\section{Introduction}
\label{sec:intro}
Online reviews shape where people stay, eat and shop, and the ratings they produce translate directly into revenue. This creates a strong incentive to manipulate them. Luca and Zervas found that about 16\% of restaurant reviews on Yelp in their Boston sample were filtered as suspicious, and that restaurants with weak or declining reputations were more likely to engage in review fraud \cite{luca2016}. He et al.\ documented an organised market in which sellers buy fake Amazon reviews through large private Facebook groups; ratings and sales rank rise while purchases continue, but after they stop, average ratings fall and the share of one-star reviews increases, which indicates that fake reviews are used mainly by low-quality products and harm consumers \cite{he2022market}. Regulators now treat the problem as consumer deception: the 2024 rule of the U.S.\ Federal Trade Commission prohibits writing, buying and selling fake consumer reviews, including reviews generated by artificial intelligence, and allows civil penalties for knowing violations \cite{ftc2024}. Enforcing such rules at platform scale requires automated detectors whose accuracy, reliability and cost are known, because the volume of incoming reviews far exceeds what human moderators can read.

Early computational work established that review spam can be learned from data, but also that the definition of the label shapes what is learned. Jindal and Liu analysed several million Amazon reviews and, lacking verified fakes, used duplicate and near-duplicate reviews as proxy labels for spam \cite{jindal2008}. Ott et al.\ removed this ambiguity by commissioning deceptive hotel reviews through Amazon Mechanical Turk and pairing them with truthful reviews; an $n$-gram classifier reached nearly 90\% accuracy on this gold-standard corpus \cite{ref22}, and the later negative-sentiment extension produced the 1,600-review Deceptive Opinion Spam corpus \cite{ott2013}. Feng et al.\ raised accuracy on the same hotel data to 91.2\% with syntactic features derived from parse trees \cite{feng2012}, and Li et al.\ built a three-domain gold standard containing customer, crowd-worker and domain-expert reviews to search for cues that generalise across domains \cite{li2014}.

A second line of work questioned whether crowdsourced fakes represent commercial fraud. Mukherjee et al.\ trained the same kind of supervised models on reviews filtered by Yelp and found that linguistic $n$-gram features, which work very well on Mechanical Turk reviews, were much less effective on real filtered reviews, whereas behavioural features of reviewers performed well \cite{mukherjee2013}. Related studies showed that spammers often act in coordinated groups \cite{mukherjee2012}, and Rayana and Akoglu combined review text, behavioural metadata and the reviewer--product network in the SpEagle framework, evaluated on Yelp collections that include the YelpChi data planned for this project \cite{rayana2015}. Together these studies show that detection performance depends on the source of the labels and on whether metadata are available, and that results obtained on one kind of fake review should not be generalised to another without testing.

Pretrained transformers changed the text side of the problem. BERT introduced deep bidirectional pretraining followed by task-specific fine-tuning \cite{devlin2019}, and RoBERTa showed that the same architecture improves substantially when pretrained longer, on more data and with dynamic masking \cite{liu2019roberta}. Fine-tuned transformers now report the highest scores on deception benchmarks, often above 90\% accuracy on the Deceptive Opinion Spam corpus (Section~\ref{sec:litreview}). On a corpus of only 1,600 short reviews, however, a well-regularised linear model over word $n$-grams is already strong, and it is not obvious that a model with about 125 million parameters adds enough accuracy to justify its training and inference cost.

This question has a big-data dimension. The difficulty for platforms is less the size of any single benchmark than the velocity of incoming text: the Boston sample of Luca and Zervas alone contained 316,415 reviews of 3,625 restaurants \cite{luca2016}, and every new review must be screened as it arrives. A detector that is marginally more accurate but two orders of magnitude slower changes the hardware and latency budget of the whole moderation pipeline, so accuracy and computational cost should be measured together. The choice of metric matters as well. Saito and Rehmsmeier showed that precision--recall analysis is more informative than ROC analysis when the positive class is rare \cite{saito2015}, which is the normal situation for fake reviews on real platforms. Metrics that average over both classes, such as macro-F1, and threshold-free summaries of the precision--recall curve, such as average precision, are therefore preferable to raw accuracy even when a benchmark happens to be balanced.

The answer also depends on evaluation choices that the field often leaves implicit. Kapoor and Narayanan found data leakage in 294 papers across 17 scientific fields and showed that, once leakage was corrected, complex machine-learning models in a reproduced case study no longer outperformed logistic regression \cite{kapoor2023}. In the Deceptive Opinion Spam corpus, every hotel contributes truthful and deceptive reviews, so a random review-level split places reviews of the same hotel in both training and test data and lets a classifier exploit hotel-specific vocabulary. Dror et al.\ surveyed ACL and TACL papers and found that statistical significance testing is often ignored or misused in natural language processing (NLP) \cite{dror2018}; a single higher score on one test set is therefore weak evidence that one model is better. Documentation practices such as datasheets for datasets make provenance and labelling decisions visible \cite{gebru2021}, which matters here because the deceptive and truthful reviews in this corpus were collected from different sources.

Held-out accuracy also says little about how a model behaves when its input changes in ways that should not matter. Ribeiro et al.\ argued that held-out accuracy often overestimates the capability of NLP models and proposed CheckList, which tests models with targeted input changes, including invariance tests in which a prediction should not change; such tests exposed many failures in commercial and state-of-the-art systems \cite{ribeiro2020}. Jin et al.\ showed that meaning-preserving word substitutions can sharply reduce the accuracy of fine-tuned BERT classifiers \cite{jin2020}. For review moderation this matters because fake-review writers control their text and can vary its surface form cheaply, while benchmark corpora such as Deceptive Opinion Spam are uniformly formatted. Whether a transformer's benchmark advantage survives simple, content-preserving changes is therefore part of the question of which model to deploy.

\emph{Research gap.} Prior work shows that transformers score highly on deceptive-review benchmarks and that linear $n$-gram models are strong on the same data \cite{ref22,feng2012}, but these results were obtained under different partitions and are usually reported without paired significance tests. Comparisons also rarely hold out entire hotels, test whether predictions survive content-preserving changes to the input, or report how much the gain costs in training time, memory and throughput. It is therefore unclear whether a fine-tuned transformer is reliably better than a fixed linear baseline when both are trained and tested on the same leakage-controlled partitions, how stable and robust any gain is, and what it costs.

\emph{Aim and research questions.} The aim of this study is to determine whether fine-tuned RoBERTa-base detects deceptive hotel reviews more accurately than a fixed TF--IDF logistic-regression baseline when both are evaluated on identical hotel-disjoint partitions, and at what computational cost. The objectives are to (i) build an audited, leakage-controlled and reproducible evaluation pipeline for the Deceptive Opinion Spam corpus; (ii) train and evaluate both models under one fixed protocol; (iii) quantify uncertainty with hotel-level resampling, paired randomisation tests and repeated partitions; and (iv) measure robustness to formatting changes and computational cost. Three research questions follow:
\begin{itemize}
\item \textbf{RQ1:} Does fine-tuned RoBERTa-base achieve higher macro-F1 than TF--IDF with logistic regression on identical hotel-disjoint test sets?
\item \textbf{RQ2:} How stable are both models across partitions, held-out hotels, review polarity and formatting perturbations?
\item \textbf{RQ3:} What computational cost, in training time, memory and inference throughput, accompanies any accuracy difference?
\end{itemize}

\emph{Contribution.} This study does not propose a new architecture. Its contribution is a controlled comparison that the cited work does not provide: (1) a linear baseline and a transformer evaluated on the same deterministic, hotel-disjoint and duplicate-audited partitions; (2) paired uncertainty estimates at the hotel level, complemented by an exploratory grouped cross-validation in which all 20 hotels are tested once; (3) robustness and per-hotel error analyses that show where the models disagree; and (4) a measured cost--accuracy trade-off with released code and split files. The planned YelpChi experiment and star-rating ablation from the original protocol could not be run because the text-containing YelpChi data were not obtained; this deviation is reported rather than concealed.

The remainder of the article is organised as follows. Section~\ref{sec:litreview} reviews recent survey and empirical work on fake-review detection. Section~\ref{sec:methods} describes the data, partitioning, models, metrics and statistical procedures. Section~\ref{sec:results} reports the results for each research question. Section~\ref{sec:conclusion} answers the research questions, discusses limitations and outlines future work, followed by the data and code availability statements.
